\section{LLM-Based Agents and Multi-Agent Systems}
\label{sec:bg-mas}

\subsection{LLM-Based Agents}
\label{sec:bg-agents}

An \emph{LLM-based agent} is a system that uses an LLM to decide how to act towards a goal. A common framework describes such an agent by four modules~\cite{wang2024survey}: a \emph{profile} that defines its role and personality, usually in the system prompt; a \emph{memory} that holds what it has observed and done; a \emph{planning} module that decomposes the goal into steps; and an \emph{action} module that turns decisions into outputs or tool calls. Generative agents showed that agents equipped with a profile and a memory of their experiences can simulate believable human behavior~\cite{park2023generative}, which is the basis of simulating stakeholders with LLMs.

\subsection{Multi-Agent Systems}
\label{sec:bg-multi-agent}

A \emph{multi-agent system} divides a task among several agents with specialized roles. Specialization keeps each prompt short and focused, lets each role use the model best suited to it, and allows agents to check each other's work. Multi-agent systems differ mainly in how agents communicate:
\begin{itemize}
    \item \textbf{Sequential (pipeline)}: each agent passes its output to the next, as in the waterfall-like phases of MetaGPT~\cite{hong2024metagpt}.
    \item \textbf{Shared workspace (blackboard)}: agents read and write common artifacts, and react to changes of the artifacts they depend on.
    \item \textbf{Conversation and debate}: agents exchange messages, for example the instructor and assistant pairs of ChatDev~\cite{qian2024chatdev} or the conversable agents of AutoGen~\cite{wu2024autogen}.
    \item \textbf{Orchestrator}: a controller decides which agent acts next.
\end{itemize}
MetaGPT encodes standard operating procedures of a software company as roles (product manager, architect, engineer) that exchange structured documents~\cite{hong2024metagpt}; ChatDev organizes software development as a chat chain between role-playing agents~\cite{qian2024chatdev}. The proposed framework combines the shared-workspace pattern, since agents communicate through versioned artifacts rather than through a message history, with an explicit graph that orchestrates the order of the agents.

\subsection{Graph-Based Agent Orchestration}
\label{sec:bg-langgraph}

LangGraph~\cite{langgraph} models an agent workflow as a \emph{state graph}:
\begin{itemize}
    \item The \textbf{state} is a typed record shared by all nodes. A node returns a partial update; a \emph{reducer} of each field defines how the update is merged, for example by replacing the value or appending to a list.
    \item A \textbf{node} is a function from the state to an update. A node may call an LLM, run plain code, or both.
    \item An \textbf{edge} names the node that runs next. A \emph{conditional edge} calls a routing function on the state to choose the next node, which allows loops and branches.
    \item A \textbf{checkpointer} saves the state after every step, so that a run can be paused, resumed after a failure from its last completed step, or inspected afterwards.
\end{itemize}
Compared with a free conversation between agents, an explicit graph makes the order of the agents, the loops and their bounds visible and testable, while the conditional edges still let the run take a different path depending on the results. The proposed workflow is a LangGraph state graph with seventeen nodes (Chapter~\ref{chap:proposed-approach}).

\subsection{Human-in-the-Loop}
\label{sec:bg-hitl}

Fully automated generation leaves errors of early stages to propagate: personas that miss a key stakeholder group lead to interviews and requirements that miss it too. A \emph{human-in-the-loop} workflow pauses at chosen points, presents an intermediate artifact to a person, and continues according to their decision: approve, edit, or reject with feedback. With a checkpointer, the pause costs nothing: the run stops, its state is saved, and it resumes from that state when the decision arrives, possibly hours later. The proposed workflow places such checkpoints after the personas, the analysis and the specification, where a human judgment has the largest effect on what follows.
